\section{Delayed transforms along a plateau path}
\label{sec:paths}

The representatives from \Cref{sec:capture} will be used on intervals of
depths chosen while estimating an actual modulation. We therefore prove
interval bounds under explicit amplitude conditions, then make the bounds
simultaneous over all snapshot representatives and dyadic list caps.
Scalar moments handle small amplitudes and static plateaus; moments that
retain the main label handle the main levels on live plateaus.

Fix a plateau $e$, with normalization $\alpha=\alpha_e$, and use uniform
probability $\mathbf U_e$ on its root.  Write $\mathcal F_j$ for the spatial
prefix through depth $j$, together with the deterministic processing data at
that prefix.  In particular, every post state, its two prospective pre-child
states, and the decision to retain its jump are $\mathcal F_j$-measurable.
The sign $r_{j+1}$ is independent and fair conditional on $\mathcal F_j$.
The multiplier $\epsilon_j$ satisfies $|\epsilon_j|\le1$ and is
$\mathcal F_{j+1}$-measurable.  It is the same for every label coordinate.
Thus a coefficient formed at depth $j$ has its output multiplied by the
later sign $r_{j+2}$.  The additional bit allowed in the finite model supplies
this sign even for the final contributing jump.

All moments below are divided by $\alpha$.  At a jumping post state the
total mass is at most $2$, and at either prospective pre-child it is at
most $4$.  Intervening removals have normalized masses $\xi_p$; full removal
at an exit can be included.  We write $K_e=\sum_p\xi_p$ for the sum along
the entire plateau path.  These quantities and the local budget statements
are those of \Cref{lemma:07-local-budget}.  No realized input membership or
noisy label is sampled on this uniform output path.
All statements use the sufficiently large $d$ convention already fixed;
their constants are uniform in the finite spatial depth and in every
integer-frequency cutoff.

\subsection{The predictable sign inequality}

We use the standard exponential-supermartingale proof of a maximal
inequality, based on Ville's stopping argument
\cite[Chapter~V, first section, Theorem~1, pp.~84--85]{Ville1939}.
The finite-depth version below also records the conditional and
predictable-masking forms required later.

\begin{lemma}[Predictable signs]
\label{lemma:paths-sign}
Let $(r_{i+1})$ be fair signs in a filtration $(\mathcal F_i)$, and let
$c_i$ be real $\mathcal F_i$-measurable coefficients.  For every $\theta>0$
and $t\ge0$, outside an event of probability at most $e^{-t}$,
\[
 \sum_{i<n}c_i r_{i+1}
 \le \frac{\theta}{2}\sum_{i<n}c_i^2+\frac{t}{\theta}
 \qquad\text{for every prefix }n.
\]
The same statement holds conditionally at a stopping-time origin and with
predictable stopping or masking of the coefficients.  For complex sums one
may apply the assertion to both signs of each real component, at a total
failure probability at most $4e^{-t}$.
\end{lemma}

\begin{proof}
The elementary inequality $\cosh x\le e^{x^2/2}$ shows that
\[
 \exp\left(\theta\sum_{i<n}c_i r_{i+1}
                  -\frac{\theta^2}{2}\sum_{i<n}c_i^2\right)
\]
is a nonnegative supermartingale starting at one.  Stop at the first
crossing of $e^t$ to obtain the asserted maximal inequality.  Everything
is in finite depth, so optional sampling needs no limiting argument.
Starting the same calculation with conditional expectation given the
origin proves the conditional assertion.  Predictable masks are absorbed
in $c_i$.
\end{proof}

\subsection{Scalar moments and their masked transforms}

\begin{proposition}[Scalar plateau paths]
\label{prop:scalar-path}
Let $0<b\le1$ be a scalar level, and let $s$ be a representative selected at
a predictable snapshot birth on a live or static plateau.  At a jumping
post state put
\[
 f_p=\frac{\mathbb E_p[a\chi_s]}{\alpha},\qquad
 u_p=\frac{f_{p,+}-f_{p,-}}2,\qquad
 \eta_p=\mathbf1_{\{\max(|f_{p,+}|,|f_{p,-}|)\le C_3b\}},
\]
where $C_3$ is a fixed constant and $f_{p,\pm}=f_p\pm u_p$ are prospective
pre-child moments.  For every $t\ge0$, with conditional failure probability
at most $Ce^{-t}$ at the birth, every contiguous interval $I$ of subsequent
plateau jumps satisfies
\begin{equation}
 \left|\sum_{p\in I}\eta_p u_p\epsilon_{j(p)}r_{j(p)+2}\right|
 \le C b^{1/4}(1+K_e+t).
 \label{eq:source-15}
\end{equation}
Constants are independent of depth, $s$, and $b$; they may depend on the
fixed mask constant $C_3$.
\end{proposition}

\begin{proof}
Identify a complex moment with a point in $\mathbb R^2$ and set
\[
 \psi(x)=b^{1/2}\big((|x|^2+b^2)^{3/4}-b^{3/2}\big).
\]
On $|x|\le4$ one has $0\le\psi(x)\le Cb^{1/2}$ and
$|\nabla\psi(x)|\le Cb^{1/2}$.  More precisely,
\[
 \nabla^2\psi(x)
 =\frac32 b^{1/2}(|x|^2+b^2)^{-5/4}
       \left((|x|^2+b^2)\operatorname{Id}-\frac12 x x^{\mathsf T}\right).
\]
Its least eigenvalue is at least
$\frac34 b^{1/2}(|x|^2+b^2)^{-1/4}$.  Consequently it is at least
$c b^{1/2}$ throughout that ball, and at least $c(C_3)>0$ on
$|x|\le C_3b$.

Define the nonnegative averaging drift and the signed fluctuation by
\[
 \delta_p=\frac{\psi(f_p+u_p)+\psi(f_p-u_p)}2-\psi(f_p),\qquad
 n_p=\frac{\psi(f_p+u_p)-\psi(f_p-u_p)}2.
\]
Taylor's formula on the two segments $f_p\pm t u_p$, $0\le t\le1$, gives
\begin{equation}
 \delta_p\ge c b^{1/2}|u_p|^2,
 \qquad \delta_p\ge c(C_3)\eta_p|u_p|^2,
 \qquad n_p^2\le Cb|u_p|^2\le Cb^{1/2}\delta_p.
 \label{eq:paths-scalar-drift}
\end{equation}
For the second inequality the masked segment lies in the ball of radius
$C_3b$, since its endpoints do.  The last inequality follows from the
gradient bound and the first inequality.

At the actual pre-child the potential is exactly
$\psi(f_p)+\delta_p+r_{j+1}n_p$.  A removal of normalized mass $\xi$ changes
the complex moment by at most $\xi$, and hence changes $\psi$ by at most
$Cb^{1/2}\xi$.  Telescoping from birth to any prefix end therefore gives
\[
 \sum\delta_p\le Cb^{1/2}(1+K_e)-\sum r_{j+1}n_p.
\]
Apply \Cref{lemma:paths-sign} to the negative fluctuation with
$\theta=\gamma b^{-1/2}$, where $\gamma>0$ is a sufficiently small numerical
constant.  By \Cref{eq:paths-scalar-drift} its quadratic term is at most
$\frac12\sum\delta_p$.  Outside probability $e^{-t}$ this proves, on every
prefix,
\[
 \sum\eta_p|u_p|^2\le Cb^{1/2}(1+K_e+t).
\]
Now $\eta_pu_p\epsilon_j$ is known in $\mathcal F_{j+1}$ before $r_{j+2}$.
Apply \Cref{lemma:paths-sign} to the real and imaginary components of this
delayed transform with $\theta=b^{-1/4}$.  Their quadratic sums are bounded
by the preceding display.  Thus every prefix transform has absolute value
at most $Cb^{1/4}(1+K_e+t)$ except on probability $Ce^{-t}$.  Subtracting
the two prefixes bounding any contiguous interval proves
\Cref{eq:source-15}.  This subtraction retains precisely the fixed mask
$\eta_p$ within the interval.
\end{proof}

\subsection{The energy of label-valued moments}

For the main levels we need an $L^1(dz)$ estimate with an absolute
coefficient on the Fisher sum. This will let the assembly charge information
over disjoint intervals while paying separately for the choice of a
representative. The label norm controls the scalar transform because
integration against a label character recovers a shifted scalar modulation
multiplied by a thermal Fourier coefficient. The main representative is
chosen so that the coefficient needed in this recovery is close to one;
\Cref{sec:assembly} gives the exact decoding and the remaining spatial
phase error. We begin with the energy identities for the label-valued moments.

Fix a main representative $s$ on a live plateau, so its main channel $A$
is fixed. At a jumping post state $p$ with post density $a_p$, define
\[
 M(z)=\frac{\mathbb E_p[a_p(Y)K_A(z-Y)]}{\alpha},\qquad
 f(z)=\frac{\mathbb E_p[a_p(Y)\chi_s(Y)K_A(z-Y)]}{\alpha}.
\]
Here the observation consists only of the main label: the rounded
coordinates used to construct snapshots have been integrated out. Define
$M_\pm,f_\pm$ by the same formulas on the two prospective children, using
the same post density $a_p$. Write
\[
 f=Mv,\qquad f_\pm=M_\pm v_\pm,\qquad |v|,|v_\pm|\le1,
 \qquad D=\frac{M_+-M_-}2,\qquad \mathfrak a=\frac{f_+-f_-}2.
\]
The functions $v$ are set to zero at zero mass.  Define
\[
 G=\int M|v|^2\,dz,\qquad J=\int\frac{D^2}{M}\,dz,
 \qquad G_\pm=\int M_\pm|v_\pm|^2\,dz.
\]
In particular $J$ agrees with \Cref{eq:source-13}.  Quotients with zero
denominator are assigned zero; nonnegativity and child averaging make
their numerators zero as well.

\begin{lemma}[Energy drift and innovations]
\label{lemma:paths-innovation}
Put
\[
 g_z=\frac12\big(M_+|v_+-v|^2+M_-|v_--v|^2\big),\qquad
 g=\int g_z\,dz,\qquad U_p=\frac{G_+-G_-}2.
\]
Then
\begin{align}
 &g_z=\frac{M_+|v_+|^2+M_-|v_-|^2}{2}-M|v|^2\ge0,
       &|\mathfrak a-vD|^2&\le M g_z,\label{eq:paths-innovation-pointwise}\\
 &G_{\mathrm{actual\ pre}}=G+g+r_{j+1}U_p,
       &|U_p|^2&\le C\max(G_+,G_-)(J+g),\label{eq:paths-energy-noise}\\
 &&\left(\int|\mathfrak a|\,dz\right)^2&\le C(GJ+g).
       \label{eq:paths-moment-variation}
\end{align}
A removal of normalized mass $\xi$ changes $G$ in absolute value by at
most $3\xi$.  From any live jumping post state,
\begin{equation}
 \mathbb E\sum_{\mathrm{future}}g\le C,
 \qquad \mathbb E\sum_{\mathrm{future}}J\le Cd,
 \qquad \mathbb E\sum_{\mathrm{future}}\xi\le2,
 \label{eq:paths-forward-budgets}
\end{equation}
where the conditional law is spatially uniform on that state.
\end{lemma}

\begin{proof}
Child averaging gives
$M_+(v_+-v)=\mathfrak a-vD$ and $M_-(v_--v)=-(\mathfrak a-vD)$.  Expansion around $v$ proves the
first identity.  If both child masses are positive, it also gives
\[
 g_z=\frac{M|\mathfrak a-vD|^2}{M^2-D^2}.
\]
This implies the innovation inequality.  If a child mass is zero, the
corresponding moment is zero and $\mathfrak a-vD=0$, giving the same inequality.
Write $e=\mathfrak a-vD$.  Expanding the difference of the two child energies yields
\[
 U_p=\int D|v|^2\,dz+2\operatorname{Re}\int\overline v e\,dz+R,
 \qquad |R|\le g.
\]
Cauchy--Schwarz bounds the first two integrals in absolute value by
$\sqrt{GJ}$ and $\sqrt{Gg}$, respectively.  Since
$G+g=(G_++G_-)/2\le\max(G_+,G_-)$, squaring proves
\Cref{eq:paths-energy-noise}.  The decomposition $\mathfrak a=vD+e$ gives
\[
 \int|\mathfrak a|\,dz\le\sqrt{GJ}+\sqrt{\left(\int M\,dz\right)g}
                  \le\sqrt{GJ}+\sqrt{2g},
\]
which proves \Cref{eq:paths-moment-variation}.

For the removal assertion, the energy has the variational expression
\[
 \int\frac{|f|^2}{M}\,dz
 =\sup_{|\gamma(z)|\le1}
       \int\big(2\operatorname{Re}(\overline\gamma f)-|\gamma|^2M\big)\,dz.
\]
Removing $k_M\ge0$ and $k_f$ with $|k_f|\le k_M$ changes every test in this
supremum by at most $3\int k_M=3\xi$.  Both the old and new moment-to-mass
ratios have modulus at most one, so the same test class represents both
energies.  Finally telescope \Cref{eq:paths-energy-noise} over the entire
future stopped plateau path.  Its martingale terms have zero conditional
expectation, its terminal pre energy is at most $4$, and the removals have
expected total at most $2$.  This gives the first bound in
\Cref{eq:paths-forward-budgets}.  The other two are
\Cref{lemma:07-local-budget}.
\end{proof}

\subsection{A vector estimate on an outer-valid run}

For a main dyadic level $b\in[b_0,1]$, declare a jump exceptional when
$b<1$ and it is bad for the \emph{main} criterion at level $2b$.  Put
\[
 e'_p=C b_0^{-2}e_{p,2b}\mathbf1_{\{p\text{ exceptional}\}}
       \quad(b<1),\qquad e'_p=0\quad(b=1).
\]
Choose $C$ so that $e'_p\ge1$ on every exceptional jump; this is possible
because a bad main jump has edge mass at least $c(2b)^2$.  Let
$E_{e,b}=\sum_p e'_p$ on the whole plateau path.  Fix a sufficiently large
numerical $C_4$.  For the specified $s,b$, call a jump \emph{inner-valid}
if it is exceptional or $\max(G_+,G_-)\le C_4b^2$, and \emph{outer-valid}
if it is exceptional or $\max(G_+,G_-)\le2C_4b^2$.  At $b=1$ choose $C_4$
large enough that every jump is inner-valid.

\begin{lemma}[Drift on a run]
\label{lemma:paths-run-drift}
Start predictably at an outer-valid jump, and stop before the first
outer-invalid jump or at plateau end.  For $t\ge0$, outside conditional
probability $e^{-t}$, every prefix $P$ of this run satisfies
\begin{equation}
 \sum_{p\in P}g_p
 \le C\left(1+\sum_{p\in P}\xi_p+\sum_{p\in P}e'_p\right)
       +Cb\sum_{p\in P}J_p+Cb t.
 \label{eq:source-16}
\end{equation}
Removals here include those between consecutive jumps and may also be
enlarged to the whole plateau sum.
\end{lemma}

\begin{proof}
Telescope the exact energy identity of
\Cref{eq:paths-energy-noise}.  Initial and terminal energies are bounded
by $4$ and removals cost at most $3\xi$.  On an exceptional jump,
$|U_p|\le4\le4e'_p$, so pay its noise directly.  On the remaining jumps,
\Cref{eq:paths-energy-noise} gives
$|U_p|^2\le Cb^2(J_p+g_p)$.  Apply \Cref{lemma:paths-sign} to minus this
remaining noise with parameter $\theta=\gamma/b$.  Its quadratic bound is
$C\gamma b\sum(J_p+g_p)$.  Choose $\gamma$ small enough to absorb the
$g$ term, using $b\le1$.  Its other term is $Cb t$, which proves the
simultaneous prefix assertion.
\end{proof}

\begin{lemma}[Delayed norm estimate on a run]
\label{lemma:paths-run-norm}
Under the hypotheses of \Cref{lemma:paths-run-drift}, set
$\lambda=b^{1/2}$ and
$W_P(z)=\sum_{p\in P}\mathfrak a_p(z)\epsilon_{j(p)}r_{j(p)+2}$.  With conditional
failure probability at most $Ce^{-t}$, simultaneously for every run prefix,
\begin{equation}
 \int|W_P(z)|\,dz
 \le C\sum_{p\in P}J_p
      +C\lambda^{-1}\left(1+\sum_{p\in P}\xi_p+\sum_{p\in P}e'_p\right)
      +C\lambda t.
 \label{eq:source-17}
\end{equation}
In particular, the coefficient of the Fisher sum is absolute, independent
of $b$ and of the number of representatives.
\end{lemma}

\begin{proof}
We prove the norm estimate using both signs of a two-step block.  Separate
contributing depths into their two parities.  For one parity start at its
first depth in the run and group successive spatial jumps into blocks
$j,j+1$, with $j$ contributing to this parity.  Only complete blocks lying
in the prefix will be used in the norm telescope.  Denote the accumulated
transform for this parity at the block start by $W^{(0)}$, and put
\[
 X_0=(W^{(0)},f_0,\lambda M_0),\qquad
 \Phi=|X_0|=\big(|W^{(0)}|^2+|f_0|^2+\lambda^2M_0^2\big)^{1/2}.
\]
The complex coordinates are viewed as pairs of real coordinates.  Subscript
$0$ denotes the first post state.  After its sign $r=r_{j+1}$, write
$(\bar M,\bar f)$ for the actual pre-child, and $(M_1,f_1)$ for its post
state after removals $(k_M,k_f)$, where
\[
 \bar M=M_0+rD_0,\quad \bar f=f_0+r\mathfrak a_0,\quad
 M_1=\bar M-k_M,\quad f_1=\bar f-k_f,
 \quad k_M\ge0,\quad |k_f|\le k_M.
\]
Here and in this proof these are pointwise identities in $z$; in particular
$\int k_M=\xi_{j+1}$.  At the end of the second jump, before its child's
processing, the change in $X_0$ is exactly
\begin{equation}
 \Delta=\big(\mathfrak a_0\epsilon_jr',\;r\mathfrak a_0-k_f+r'\mathfrak a_1,
                    \lambda(rD_0-k_M+r'D_1)\big),\qquad r'=r_{j+2}.
 \label{eq:paths-two-step-increment}
\end{equation}

We first prove the required pointwise quadratic estimate.  Put
\[
 q=(|v_0|^2+\lambda^2)^{1/2},\qquad H=M_0q\le\Phi,
 \qquad j_i=\frac{D_i^2}{M_i},\qquad e_i=\mathfrak a_i-v_iD_i.
\]
For $M_0>0$, $q\ge\lambda$ and $|v_0|^2/q\le1$.  Positivity of the child
masses and of the removed density gives
\begin{equation}
 0\le M_1\le\bar M\le2M_0,
 \quad 0\le k_M\le\bar M\le2M_0,
 \quad |D_i|\le M_i,
 \quad |e_i|^2\le M_i g_{i,z}.
 \label{eq:paths-positive-domination}
\end{equation}
There is also a comparison of the child's post mean with the starting
mean:
\begin{equation}
                 M_1|v_1-v_0|^2\le4g_{0,z}+8k_M.
 \label{eq:paths-child-mean}
\end{equation}
To verify it, let $\bar v=\bar f/\bar M$ when $\bar M>0$, and zero otherwise.
The actual
child's term in the definition of $g_{0,z}$ gives
$\bar M|\bar v-v_0|^2\le2g_{0,z}$.  If $M_1=0$, the left side of
\Cref{eq:paths-child-mean} is zero.  If $0<M_1\le\bar M/2$, then
$M_1\le k_M$ and the bound $|v_1-\bar v|\le2$ gives
$M_1|v_1-\bar v|^2\le4k_M$.  If $M_1>\bar M/2$, the identity
\[
 M_1(v_1-\bar v)=k_M\bar v-k_f
\]
instead gives $M_1|v_1-\bar v|^2\le4k_M^2/M_1\le4k_M$, since $k_M<M_1$.
Thus the same bound covers every intermediate or arbitrarily small
remaining mass.  The squared triangle inequality, and $M_1\le\bar M$,
now prove \Cref{eq:paths-child-mean}.

Splitting $\mathfrak a_0=v_0D_0+e_0$ and retaining the factor $q$ yields
\begin{equation}
 \frac{|\mathfrak a_0|^2}{H}\le2j_0+\frac{2g_{0,z}}\lambda.
 \label{eq:paths-first-remainder}
\end{equation}
For the child jump use the three terms
$\mathfrak a_1=v_0D_1+(v_1-v_0)D_1+e_1$.  Their separate squared ratios obey
\begin{align*}
 \frac{|v_0D_1|^2}{H}&\le2j_1,\\
 \frac{|v_1-v_0|^2D_1^2}{H}
   &\le\frac{2}{\lambda}M_1|v_1-v_0|^2
       \le\frac{8g_{0,z}+16k_M}{\lambda},\\
 \frac{|e_1|^2}{H}&\le\frac{2g_{1,z}}\lambda.
\end{align*}
The factor $M_1/M_0\le2$ supplies the first and third bounds; the second
uses $|D_1|\le M_1$ and \Cref{eq:paths-child-mean}.  In particular no lower
bound on $M_1/M_0$ is used.  Hence
\begin{equation}
 \frac{|\mathfrak a_1|^2}{H}
 \le6j_1+\frac{24g_{0,z}+6g_{1,z}+48k_M}{\lambda}.
 \label{eq:paths-second-remainder}
\end{equation}
The mass-coordinate and removal terms satisfy
\[
 \frac{\lambda^2D_0^2}{H}\le j_0,
 \qquad \frac{\lambda^2D_1^2}{H}\le2j_1,
 \qquad \frac{k_M^2+|k_f|^2}{H}\le\frac{4k_M}{\lambda}.
\]
These follow again from \Cref{eq:paths-positive-domination} and
$0<\lambda\le1$.  If $M_0=0$, all the masses, half-differences, and
removals in \Cref{eq:paths-two-step-increment} are zero by positivity, so
$\Delta=0$ whether or not $W^{(0)}$ is zero.  This deals with every zero
denominator without an approximation argument.

For $\Phi>0$, concavity of the square root gives the global norm bound
\[
 |X_0+\Delta|\le\Phi+\left\langle\frac{X_0}{\Phi},\Delta\right\rangle
                              +\frac{|\Delta|^2}{2\Phi}.
\]
Combining the preceding pointwise inequalities and integrating proves
\begin{equation}
 \int\frac{|\Delta|^2}{2\Phi}\,dz
 \le C(J_0+J_1)+C\lambda^{-1}(g_0+g_1+\xi_{j+1}).
 \label{eq:paths-norm-remainder}
\end{equation}
The absolute Fisher coefficient results from the estimates of $v_0D_i$
with $q$, rather than replacing $q$ by $\lambda$ in those terms.

It remains to control the integrated linear term and its timing.  Write
$\Gamma_0=X_0/\Phi$, with value zero where $\Phi=0$.  Its norm is at most
one.  The linear term is
\[
 r L_0+r'L_1+L_{\rm kill},
\]
where $L_{\rm kill}$ has absolute value at most $2\xi_{j+1}$, and
$L_0,L_1$ are real integrated coefficients.  Specifically $L_0$ pairs
$\Gamma_0$ with $(0,\mathfrak a_0,\lambda D_0)$, whereas $L_1$ pairs it with
$(\mathfrak a_0\epsilon_j,\mathfrak a_1,\lambda D_1)$.  By
\Cref{eq:paths-moment-variation} and
$(\int|D_i|)^2\le(\int M_i)J_i\le2J_i$, their squares satisfy
\begin{equation}
 |L_0|^2+|L_1|^2
 \le C\sum_{i=0,1}\big((\lambda^2+G_i)J_i+g_i\big).
 \label{eq:paths-linear-variance}
\end{equation}

At a block start $W^{(0)}$ is $\mathcal F_j$-measurable: the previous
output of the same parity, if present, used $r_j$.  Thus $L_0$ is known
before $r_{j+1}$.  After that bit the removals, second post state, and
$\epsilon_j$ are all known, so $L_1$ is known before $r_{j+2}$.  Form one
scalar sign scheme by including $rL_0$ whenever the first jump is in the
run, and including $r'L_1$ only if the second jump is also in the run.
Both inclusion decisions are known before their respective signs.  The
linear sum through the last complete block of any prefix is itself a
prefix of this scheme.  If the run stops after a first step, no first
coefficient is retrospectively conditioned on the next validity decision.

On a nonexceptional included state $G_i\le2C_4b^2$.  At an exceptional
state both $G_i$ and $J_i$ are at most $2$, and $e'_i\ge1$.  Therefore
\[
 G_iJ_i\le Cb^2J_i+Ce'_i.
\]
Apply \Cref{lemma:paths-sign} to the linear scheme with parameter
$1/\lambda$.  Its quadratic term, by
\Cref{eq:paths-linear-variance}, is at most
\[
 C\sum J_p+C\lambda^{-1}\sum(g_p+e'_p),
\]
because $\lambda=b^{1/2}$ and $b\le1$.  The tail term is $\lambda t$.
All sums here can be restricted to the complete blocks of the prefix.
At the first block $W^{(0)}=0$ and $\int\Phi\le C$.  Between complete
blocks, passing from a pre state to the next post state changes the
integrated norm by at most $2\xi$, by the norm's Lipschitz property.
Telescoping, using \Cref{eq:paths-norm-remainder}, proves for the completed
blocks the bound
\[
 C\sum J_p+C\lambda^{-1}\left(1+\sum\xi_p+\sum e'_p+\sum g_p\right)
                              +C\lambda t.
\]
At most one contributing output of this parity is unmatched at a prefix
end.  It costs at most $\int|\mathfrak a_p|\le\int M_p\le2$, even if its delayed
sign occurs after the run has stopped.  Add this output directly.
The other parity is started conditionally on reaching its first jump,
one step later if necessary, and obeys the same bound with the same
conditional failure estimate.  Each prefix jump is used in at most two
of the completed-block estimates.  Finally apply
\Cref{eq:source-16} on this prefix.  Its $g$ contribution becomes
$C\lambda^{-1}(1+\sum\xi+\sum e')+C(b/\lambda)\sum J+C(b/\lambda)t$.
Since $b/\lambda=\lambda\le1$, adding the two parity bounds gives
\Cref{eq:source-17}.
\end{proof}

\subsection{Predictable runs and Fisher checkpoints}

The run estimate charges the Fisher sum from its origin. Subtracting two
prefixes to estimate a later interval would therefore charge the earlier
part of the run twice. We introduce predictable checkpoints so that this
extra Fisher charge is bounded by an absolute constant.

\begin{proposition}[Vector estimates on arbitrary inner-valid intervals]
\label{prop:vector-path}
Fix a live plateau, a main level $b\in[b_0,1]$, and a representative $s$
at its snapshot birth.  There is a family of predictable estimate origins,
containing run starts and the checkpoints described below, whose
conditional expected size at the birth is at most
$P(d,b_0^{-1})$ for an absolute fixed polynomial $P$.
For every $t\ge0$, outside a conditional event of probability at most
$C P(d,b_0^{-1})e^{-t}$, every contiguous interval $I$ of inner-valid jumps
after this birth satisfies
\begin{equation}
 \int\left|\sum_{p\in I}\mathfrak a_p(z)\epsilon_{j(p)}r_{j(p)+2}\right|\,dz
 \le C\sum_{p\in I}J_p
       +Cb^{-1/2}(1+K_e+E_{e,b})+Cb^{1/2}t.
 \label{eq:paths-inner-interval}
\end{equation}
The Fisher sum is taken only over $I$.  The representative remains
eligible for these estimates until the plateau ends, even if another
snapshot has meanwhile replaced the one at which it was born.
\end{proposition}

\begin{proof}
Start a run at the first inner-valid opportunity after birth, possibly at
the birth jump itself.  Follow consecutive outer-valid jumps, stopping
before the first outer-invalid jump or at plateau end.  After a break,
start the next run at the first later inner-valid opportunity.  These are
predictable rules: both prospective child energies and the exceptional
indicator are known before the jump sign.  Every contiguous inner-valid
interval lies in one of these outer-valid runs.

We first bound the number of run starts conditionally on the birth
prefix.  A start on an exceptional jump can be charged to that jump.
\Cref{lemma:07-edge-charges} bounds the expected number of these by
$d\poly(h)$; there are none when $b=1$.  At a nonexceptional start,
$G_{\rm start}\le C_4b^2$.  If its run breaks, include the first invalid
jump in a comparison interval, although the transform estimate stops
before it.  At that jump $\max(G_+,G_-)>2C_4b^2$.  Conditional on the
prefix there, with probability at least $1/2$ its actual pre-child energy
exceeds $2C_4b^2$.  Mark such a break.  The comparison intervals for
different runs are disjoint, including their invalid jumps.

Let $Q$ denote one such start-to-break comparison interval, and put
$\mathcal M_Q=\sum_{p\in Q}r_{j(p)+1}U_p$.  By the exact energy identity
and the removal bound, a marked break implies
\[
 C_4b^2\le\sum_{p\in Q}g_p+3\sum_{p\in Q}\xi_p+|\mathcal M_Q|.
\]
Consequently, with a numerical constant depending only on $C_4$,
\begin{equation}
 \mathbf1_{\{Q\text{ marked}\}}
 \le Cb^{-2}\sum_{p\in Q}(g_p+\xi_p)+Cb^{-4}|\mathcal M_Q|^2.
 \label{eq:paths-marked-break}
\end{equation}
To use orthogonality, first include every start-to-break interval of this
type and the possible unfinished terminal interval, without conditioning
on a mark.  Its starting rule and inclusion through the invalid bit are
predictable.  Thus stopped martingale orthogonality gives
\[
 \mathbb E\sum_Q|\mathcal M_Q|^2
 =\mathbb E\sum_{p\in\bigcup Q}U_p^2
 \le C\mathbb E\sum_p(J_p+g_p),
\]
using \Cref{eq:paths-energy-noise} and $\max(G_+,G_-)\le4$.  All
expectations here are conditional at birth.  Each nonexceptional break
is marked with conditional probability at least $1/2$, so
\Cref{eq:paths-marked-break,eq:paths-forward-budgets} imply
\[
 \mathbb E\#\{\text{nonexceptional-start runs that break}\}
 \le Cb^{-2}+Cdb^{-4}.
\]
At most one final run fails to break.  Including exceptional starts
proves a bound $P(d,b_0^{-1})$ for all starts, uniformly in $s$.  The
factor $\poly(h)$ is absorbed in a fixed power of $b_0^{-1}$.

Within each outer-valid run, also start an estimate whenever the
accumulated $J$ since the preceding origin reaches one.  More precisely,
after including the jump that reaches this threshold, reset the counter
and use the next jump as a new origin if it remains in the run.  Each
such origin's estimate continues to that run's same break or end.  The
counter decision is known at this next origin; no future bit is used.
The number of additional origins is at most the total $J$ of the path,
and $J_p\le\int M_p\le2$ bounds a threshold overshoot.  Therefore their
conditional expected number is at most $Cd$ by
\Cref{eq:paths-forward-budgets}.

Apply \Cref{lemma:paths-run-norm} conditionally at every origin, with the
same parameter $t$.  Summing its conditional failure probabilities
weights each origin by its probability of being reached.  The expected
origin count just proved bounds the resulting probability by
$C P(d,b_0^{-1})e^{-t}$.
On the event of simultaneous success, fix any desired interval $I$ and
choose the latest origin $q$ in its outer-valid run at or before its first
jump.  The total $J$ between $q$ and that first jump is at most $3$;
indeed the counter is below one before a threshold crossing, and one
jump adds at most two.  The interval transform is the difference of
two prefixes from $q$.  Their Fisher sums add to
\[
 \sum_{p\in I}J_p+2\sum_{q\le p<\min I}J_p
 \le\sum_{p\in I}J_p+6.
\]
Their removal and exceptional sums are bounded by twice $K_e$ and
$E_{e,b}$, respectively.  The extra constant is absorbed by
$b^{-1/2}(1+K_e+E_{e,b})$.  The triangle inequality in $L^1(dz)$ then
proves \Cref{eq:paths-inner-interval}.  This argument also permits the
interval to begin long after the snapshot's replacement.
\end{proof}

\subsection{Simultaneous representatives and random list caps}

\begin{lemma}[Parameters for all representatives at a level]
\label{lemma:paths-simultaneous}
On each plateau and at each permitted level there are nonnegative random
parameters $t_b,H_b$, independent of the actual frequency later being
estimated, with the following properties.  Almost surely the scalar
bound of \Cref{prop:scalar-path}, or the vector interval bound of
\Cref{prop:vector-path}, holds with $t=t_b$ for every representative in
the snapshot construction used on the path.  The activation conclusion
of \Cref{lemma:07-activation} holds with miss-block length $H_b$.
Moreover
\begin{align}
 \mathbb E_{\mathbf U_e}t_b
   &\le Ck_{b/4}+C\log(Cd/b) &&\text{for main levels},
       \label{eq:paths-main-parameter}\\
 \mathbb E_{\mathbf U_e}t_b
   &\le C\big(d_1+\log(1/b)\big) &&\text{for scalar levels},
       \label{eq:paths-scalar-parameter}\\
 \mathbb E_{\mathbf U_e}H_b
   &\le C\log(Cd_1/b)+C\log(1/\sigma)
       &&\text{for either criterion}.
       \label{eq:paths-miss-parameter}
\end{align}
These assertions hold simultaneously over the countably many levels.
\end{lemma}

\begin{proof}
Begin with a deterministic dyadic cap $D_{\max}=2^m$, $m\ge0$.
The broad list is constructed independently of this cap.  The capped
snapshot scheme creates snapshots only while its list length is at most
$D_{\max}$; estimates for representatives already born continue to the
plateau end.  By \Cref{lemma:07-snapshots}, the expected number of births
is bounded by a fixed polynomial in $(D_{\max},d_1,1/b)$.  The number of
representatives in each snapshot is at most
\[
 R_b^{\rm main}=\lceil e^{k_{b/2}}\rceil,
 \qquad R_b^{\rm sc}=\exp\big(C(d_1+\log(1/b))\big),
\]
as established in that lemma.  Scalar estimates originate at these
births.  For each main representative, \Cref{prop:vector-path} bounds
the conditional expected number of run and checkpoint origins by a
fixed polynomial in $(d,b_0^{-1})$.  For sufficiently large $d$ this
factor is also bounded by a fixed power of $d$, because
$b_0^{-1}\le2h^{16}\le d$.

Here is the conditional union argument, which also applies when the
number of origins is random.  Enumerate predictable origins by their
finite tree addresses and, within an address, by their representative
indices.  If $B_o$ is the event that origin $o$ is reached and an estimate
starting there fails with deterministic parameter $t$, then
\[
 \mathbb P(B_o)
 \le Ce^{-t}\mathbb P(o\text{ is reached}).
\]
Sum this inequality over origins and then use their expected count.
Thus the failure probability for a capped scheme is at most $Ce^{-t}$
times its stated representative cardinality and polynomial count.  No
union over the signed span of the broad list is taken.

Choose sufficiently large fixed constants and define baseline parameters
\[
 T_b^{\rm main}=Ck_{b/4}+C\log(Cd/b),\qquad
 T_b^{\rm sc}=C(d_1+\log(1/b)).
\]
They absorb the logarithms of the representative cardinalities and all
polynomial factors except the cap.  For a sufficiently large further
fixed $C_{\rm cap}$, use
\[
 t=T_b+C_{\rm cap}\log(CD_{\max})+z
\]
in each deterministic capped scheme.  Its failure probability is at most
$C2^{-2m}e^{-z}$, after increasing $C_{\rm cap}$ if necessary.  Summing
over all dyadic caps gives $Ce^{-z}$.  These events are nested as $z$
increases, since every right-hand side in the path estimates increases.
Let $Z_b$ be the least nonnegative integer $z$ for which all these capped
schemes succeed.  Integer tail summation gives
\[
 \mathbb P(Z_b>n)\le Ce^{-n},\qquad \mathbb E Z_b\le C,
\]
and $Z_b$ is finite almost surely.

Only now choose a cap from the realized output path: let
$\widehat D$ be the least dyadic integer at least
$1+|D|_{\rm final}$ for the broad list at this level.  Thus
$1+|D|_{\rm final}\le\widehat D<2(1+|D|_{\rm final})$, and its capped
scheme never breaches the cap on this path.  The simultaneous assertion
over deterministic caps permits this selection without a stopping-time
claim about $\widehat D$.  Set
\[
 t_b=T_b+C_{\rm cap}\log(C\widehat D)+Z_b.
\]
Broad capture gives $\mathbb E|D|_{\rm final}\le Cd_1b^{-C}$, so Jensen's
inequality yields
\[
 \mathbb E\log(C\widehat D)
 \le C+\log\big(C(1+\mathbb E|D|_{\rm final})\big)
 \le C\log(Cd_1/b).
\]
This proves \Cref{eq:paths-main-parameter,eq:paths-scalar-parameter}, using
$d_1=d$ for main levels and $\log d_1\le d_1$ for scalar levels.
Use \Cref{lemma:07-activation} with this same cap and put
\[
 H_b=C\log(C\widehat Dd_1/b)+C\log(1/\sigma).
\]
Its expectation satisfies \Cref{eq:paths-miss-parameter}.  In the static
case the optional term $\log(1/\sigma)=40\log h$ is at most $Cd\le Cd_1$;
thus it causes no loss in the scalar plateau bound used below.

For each level the construction has a probability-one success event.
Their countable intersection still has probability one.  The resulting
parameters depend on the plateau path and the level, but not on the
integer frequency whose transform will be compared to these
representatives.
\end{proof}
